\clearpage
\section[Conclusão]{Conclusão}

Como Ages I, não tenho muito parâmetro para dizer se o projeto foi ou não o
melhor trabalhado, já que foi o único até o momento. No entanto, posso garantir que
fiquei satisfeito com meu comprometimento e minhas soft skills como Ages I, sendo
capaz de ultrapassar a barreira da comunicação com integrantes de grupo, que
eram desconhecidos em sua maioria até o começo do projeto.

Achei fundamentais para o desenvolvimento, disciplinas como Gerenciamento
de Configuração de Software e uma parte de \ac{poo}, onde tivemos uma prévia de
programação funcional, foram conhecimentos importantes para chegar com uma
base decente para a Ages, ainda assim senti um pouco de falta de uma boa
introdução à programação funcional, invés de apenas uma prévia. Além disso
conceitos de Engenharia de Requisitos e Introdução à Engenharia de Software dominaram a
parte da metodologia de desenvolvimento e organização da equipe.

Pensando para o futuro, acredito que posso fazer melhor na questão da
orientação dos Ages 1, pois vi integrante do grupo Ages 1 um pouco mais tímidos
que ficaram perdidos, mas com medo de perguntar, atrapalhando parcialmente o
desenvolvimento. Além disso identifiquei que muitos componentes com tamanho
forçado representaram um problema para o final do projeto, visto que eram muito
difíceis de ser reutilizados já que muitas vezes não encaixaram na página como
esperado.

Como Ages 1 não consegui ajudar tanto os outros e nem mudar os
componentes, por medo de estragar outras páginas. Ainda assim, não deixei de
notar a complicação que estes componentes trouxeram, forçando a tomada de
decisões que não seriam ideais, mas eram as mais viáveis para o fim do projeto.

Também acredito que a falta de conhecimento do backend, me forçou a
depender da comunicação com os outros para entender como fazer a integração,
visto que não sabia ler e entender quais os parâmetros esperados. Sendo assim,
creio que um maior entendimento da lógica das ferramentas utilizadas, indiretamente
seria uma ótima melhora, visto que facilitaria a parte da integração.

Visando estes pontos a melhorar, me comprometo a tomar como ponto de
ação, ajudar os Ages 1 sempre que possível, tomando a iniciativa sempre e
trabalhar mais no backend para garantir um conhecimento de fullstack e ter mais
entendimento do desenvolvimento geral do projeto e não apenas da parte em que
estava desenvolvendo.
